% =====================================================================
%  LAB 1 -- FIR und IIR Filter Design
%  Headings copied word for word from DSP2_Lab1.pdf
% =====================================================================
\clearpage
\section{Lab exercise 1: FIR und IIR Filter Design}

% ---------------------------------------------------------------------
% Like in the DSP1 diary, each chapter starts with your own theory
% introduction (about one page) before the first subsection.
% ---------------------------------------------------------------------
TODO: chapter introduction -- theory about FIR and IIR filters, window
method, equiripple method, that will be followed in the laboratory work.

\subsection{Laboratory 1 -- The task is to compare FIR Filters with different windows and to implement an own IIR filter function.}

\subsubsection{(a) Use the functions developed in the last semester to analyze the audio file ps.wav to determine which noise signal is present.}

TODO

\subsubsection{(b) Experiment with the fdatool (rectangle window) and create a FIR filter that is able to remove the interference from the audio file.}

TODO

\subsubsection{(c) Export the filter using the File/Export menu and filter the audio file. Analyze the result visually in the frequency spectrum as well as aural.}

TODO

\subsubsection{(d) Use different windows for designing the FIR filter and compare the results and explain the advantages and disadvantages. Compare the different filters with the fvtool or with the freqz function}

TODO

\subsubsection{(e) Construct the FIR filter with the Equiripple method. Analyze the differences to the Window method.}

TODO

\subsubsection{(f) Implement your own FIR and IIR filtering function for vector-signals.}

% formula given on the exercise sheet:
\begin{equation}
  y[n] = \sum_{k=1}^{M} b_k \cdot x[n-k] - \sum_{k=1}^{M} a_k \cdot y[n-k]
  \tag{1}
\end{equation}

TODO

\subsubsection{(g) Design with the fdatool an IIR filter which can be used instead of the FIR filter. Export the IIR filter as second-order-sections (SOS).}

TODO

\subsubsection{(h) Implement your own filter function to be able to filter the audio file by using the second-order-sections. Compare the result to the sosfilt function of matlab.}

TODO

\subsection{Homework 1 (Implement your own Overlap-Add FIR-Filter function)}

\subsubsection{Implement your own Overlap-Add FIR-Filter function. Use also your own functions implemented in the winter semester. Compare it with the fftfilt function of matlab.}

TODO
